% GENERATED FILE. Edit canonical lessons, then run npm run generate:books.
% canonical-source-hash: fnv1a64:6a5ad693f3d57408
% canonical-lessons: LA-C542-strigilis, LA-C542-unguo, LA-C542-ambulatio, LA-C542-rusticor, LA-C542-amoenus

\chapter{Oil, walks and the country}
\label{ch:la-oil-walks-and-the-country}
\hlchaptermodality{542}

\begin{chapteropening}
\textbf{By the end of this chapter:} \emph{I can describe what Romans did after the baths.}

Complete the last lesson of chapter 542: I can describe what Romans did after the baths.
\end{chapteropening}

\section[strigilis, strigilis]{strigilis, strigilis — a strigil, scraper}

\label{lesson:LA-C542-strigilis}



\begin{quote}
Before the new one: say the Latin for the cold room, then the Latin for a pool.
\end{quote}

\subsection*{You'll want to know: strigilis, strigilis}
\textbf{strigilis, strigilis} — \textquotedblleft{}a strigil, scraper\textquotedblright{}.

\begin{quote}
Servus dorsum dominī strigile pūrgat.
\end{quote}

\begin{quote}
\emph{The slave cleans his master's back with a scraper.}
\end{quote}

Feminine, third declension. After oil and sweat, bathers scraped their skin clean with this curved tool.

\begin{cousinweb}
\textbf{strigilis} is built, Lewis \& Short says, on \emph{stringō} \textquotedblleft{}to strip, graze\textquotedblright{}; EtymDB compares Greek \emph{stlengis}, so the exact history is not settled.
\end{cousinweb}

\begin{grammarlens}[title={What you've built}]
One more word for telling, explaining and describing at B1.
\end{grammarlens}

\subsection*{Your turn}
\begin{itemize}
\raggedright
  \item \emph{Say it:} \emph{strigilis, strigilis}
  \item \emph{Say it:} \emph{strigilis, strigilis}, once more
  \item \emph{Say it:} a sentence of your own with \emph{strigilis, strigilis}
  \item \emph{Recall:} say the Latin for the cold room, then the Latin for a pool, then say \emph{strigilis, strigilis} again
\end{itemize}

\begin{tcolorbox}[breakable,colback=teal!4,colframe=teal!35!black,title={Before you move on}]
What does strigilis, strigilis mean? (\textquotedblleft{}A strigil, scraper\textquotedblright{}.) Say it once more.
\end{tcolorbox}
\section[unguō, unguere]{unguō, unguere — to anoint, rub with oil}

\label{lesson:LA-C542-unguo}



\begin{quote}
Before the new one: say the Latin for a pool, then the Latin for a strigil.
\end{quote}

\subsection*{You'll want to know: unguō, unguere}
\textbf{unguō, unguere} — \textquotedblleft{}to anoint, rub with oil\textquotedblright{}.

\begin{quote}
Post balneum corpus oleō unguimus.
\end{quote}

\begin{quote}
\emph{After the bath we rub the body with oil.}
\end{quote}

\emph{unguō, unguere, ūnxī, ūnctum}; also spelt \emph{ungō}. \emph{Unguentum}, already known, comes from it. English \emph{unction}.

\begin{cousinweb}
\textbf{unguō} is inherited from a Proto-Indo-European root meaning \textquotedblleft{}to anoint, salve\textquotedblright{} (EtymDB). \emph{unguentum} \textquotedblleft{}ointment\textquotedblright{} comes from it, and English \emph{unction} and \emph{ointment} from its family.
\end{cousinweb}

\begin{grammarlens}[title={What you've built}]
One more word for telling, explaining and describing at B1.
\end{grammarlens}

\subsection*{Your turn}
\begin{itemize}
\raggedright
  \item \emph{Say it:} \emph{unguō, unguere}
  \item \emph{Say it:} \emph{unguō, unguere}, once more
  \item \emph{Say it:} a sentence of your own with \emph{unguō, unguere}
  \item \emph{Recall:} say the Latin for a pool, then the Latin for a strigil, then say \emph{unguō, unguere} again
\end{itemize}

\begin{tcolorbox}[breakable,colback=teal!4,colframe=teal!35!black,title={Before you move on}]
What does unguō, unguere mean? (\textquotedblleft{}To anoint, rub with oil\textquotedblright{}.) Say it once more.
\end{tcolorbox}
\section[ambulātiō, ambulātiōnis]{ambulātiō, ambulātiōnis — a walk, stroll}

\label{lesson:LA-C542-ambulatio}



\begin{quote}
Before the new one: say the Latin for a strigil, then the Latin for to anoint.
\end{quote}

\subsection*{You'll want to know: ambulātiō, ambulātiōnis}
\textbf{ambulātiō, ambulātiōnis} — \textquotedblleft{}a walk, stroll\textquotedblright{}.

\begin{quote}
Post cēnam ambulātiōnem in hortō facimus.
\end{quote}

\begin{quote}
\emph{After dinner we take a walk in the garden.}
\end{quote}

Feminine, third declension, from \emph{ambulō}. Also a covered walk in a villa.

\begin{cousinweb}
\textbf{ambulātiō} is the action noun of \emph{ambulō} \textquotedblleft{}to walk\textquotedblright{} (Lewis \& Short; EtymDB). English \emph{ambulance} and \emph{perambulate} come from the same verb.
\end{cousinweb}

\begin{grammarlens}[title={What you've built}]
One more word for telling, explaining and describing at B1.
\end{grammarlens}

\subsection*{Your turn}
\begin{itemize}
\raggedright
  \item \emph{Say it:} \emph{ambulātiō, ambulātiōnis}
  \item \emph{Say it:} \emph{ambulātiō, ambulātiōnis}, once more
  \item \emph{Say it:} a sentence of your own with \emph{ambulātiō, ambulātiōnis}
  \item \emph{Recall:} say the Latin for a strigil, then the Latin for to anoint, then say \emph{ambulātiō, ambulātiōnis} again
\end{itemize}

\begin{tcolorbox}[breakable,colback=teal!4,colframe=teal!35!black,title={Before you move on}]
What does ambulātiō, ambulātiōnis mean? (\textquotedblleft{}A walk, stroll\textquotedblright{}.) Say it once more.
\end{tcolorbox}
\section[rūsticor, rūsticārī]{rūsticor, rūsticārī — to stay in the country}

\label{lesson:LA-C542-rusticor}



\begin{quote}
Before the new one: say the Latin for to anoint, then the Latin for a walk.
\end{quote}

\subsection*{You'll want to know: rūsticor, rūsticārī}
\textbf{rūsticor, rūsticārī} — \textquotedblleft{}to stay in the country\textquotedblright{}.

\begin{quote}
Aestāte pater in vīllā rūsticātur.
\end{quote}

\begin{quote}
\emph{In summer father stays in the country at the villa.}
\end{quote}

Deponent: \emph{rūsticor, rūsticārī, rūsticātus sum}, from \emph{rūsticus}. English \emph{rusticate}.

\begin{cousinweb}
\textbf{rūsticor} is built on \emph{rūsticus} \textquotedblleft{}of the country\textquotedblright{}, from \emph{rūs} \textquotedblleft{}countryside\textquotedblright{} (Lewis \& Short). English \emph{rusticate} comes from it.
\end{cousinweb}

\begin{grammarlens}[title={What you've built}]
One more word for telling, explaining and describing at B1.
\end{grammarlens}

\subsection*{Your turn}
\begin{itemize}
\raggedright
  \item \emph{Say it:} \emph{rūsticor, rūsticārī}
  \item \emph{Say it:} \emph{rūsticor, rūsticārī}, once more
  \item \emph{Say it:} a sentence of your own with \emph{rūsticor, rūsticārī}
  \item \emph{Recall:} say the Latin for to anoint, then the Latin for a walk, then say \emph{rūsticor, rūsticārī} again
\end{itemize}

\begin{tcolorbox}[breakable,colback=teal!4,colframe=teal!35!black,title={Before you move on}]
What does rūsticor, rūsticārī mean? (\textquotedblleft{}To stay in the country\textquotedblright{}.) Say it once more.
\end{tcolorbox}
\section[amoenus, amoena, amoenum]{amoenus, amoena, amoenum — lovely, pleasant (of places)}

\label{lesson:LA-C542-amoenus}



\begin{quote}
Before the new one: say the Latin for a walk, then the Latin for to stay in the country.
\end{quote}

\subsection*{You'll want to know: amoenus, amoena, amoenum}
\textbf{amoenus, amoena, amoenum} — \textquotedblleft{}lovely, pleasant (of places)\textquotedblright{}.

\begin{quote}
Prope rīvum locus amoenus est.
\end{quote}

\begin{quote}
\emph{There is a lovely spot near the stream.}
\end{quote}

Of gardens, groves and views. \emph{Locus amoenus} — a pleasant place.

\begin{cousinweb}
\textbf{amoenus} is a native Latin adjective; Lewis \& Short reports links with \emph{amō} \textquotedblleft{}to love\textquotedblright{} and with a Greek word, but the origin is not settled.
\end{cousinweb}

\begin{grammarlens}[title={What you've built}]
One more word for telling, explaining and describing at B1.
\end{grammarlens}

\subsection*{Your turn}
\begin{itemize}
\raggedright
  \item \emph{Say it:} \emph{amoenus, amoena, amoenum}
  \item \emph{Say it:} \emph{amoenus, amoena, amoenum}, once more
  \item \emph{Say it:} a sentence of your own with \emph{amoenus, amoena, amoenum}
  \item \emph{Recall:} say the Latin for a walk, then the Latin for to stay in the country, then say \emph{amoenus, amoena, amoenum} again
\end{itemize}

\begin{tcolorbox}[breakable,colback=teal!4,colframe=teal!35!black,title={Before you move on}]
What does amoenus, amoena, amoenum mean? (\textquotedblleft{}Lovely, pleasant (of places)\textquotedblright{}.) Say it once more.
\end{tcolorbox}
